\begin{itemize}
\item \textbf{Komutační křivka:} spojnice vrcholů ustálených souměrných hysterezních smyček získaných při různých amplitudách pole. Téměř se shoduje s křivkou prvotní magnetizace, ale snáze a přesněji se měří, proto ji v praxi zastupuje.
\item \textbf{Dynamická hysterezní smyčka:} smyčka při rychlých změnách střídavého pole. S frekvencí se oproti statické smyčce rozšiřuje, její špičky se zaoblují a roste koercitivní pole; závisí i na časovém průběhu indukce. Používá se k posouzení magnetování a ztrát při střídavém provozu.
\end{itemize}
